\documentclass{beamer}

\usetheme{Madrid}
\usecolortheme{whale}
\setbeamertemplate{navigation symbols}{}
\setbeamertemplate{caption}[numbered]

\usepackage{ctex}
\usepackage{graphicx}
\usepackage{amsmath, amssymb, amsthm}
\usepackage{booktabs}
\usepackage{hyperref}

\makeatletter
\providecommand{\framesettitle}[1]{%
  \beamer@frametitle{#1}%
}
\makeatother

\title[GDN图卷积自注意力]{基于GDN图卷积的自注意力机制研究}
\subtitle{学士学位论文答辩}
\author{王惟是}
\institute{北京师范大学文理学院}
\date{\today}

\begin{document}

\begin{frame}
  \titlepage
\end{frame}

\begin{frame}
  \tableofcontents
\end{frame}

\section{研究背景与意义}

\begin{frame}
  \framesettitle{图结构数据的普遍性}

  \begin{block}{图结构数据的普遍性}
    图结构数据在现实世界中普遍存在，图神经网络（GNN）已成为处理这类数据的主流方法。传统的卷积神经网络难以直接应用于非欧几里得空间数据，而GNN通过迭代聚合邻居节点信息，有效捕捉图的拓扑结构特征。近年来，GNN在社交网络分析、推荐系统、生物信息学等领域取得了显著成果。
  \end{block}

  \vspace{0.3cm}

  \begin{block}{典型应用场景}
    \begin{itemize}
      \item \textbf{社交网络}：用户之间的关系网络，节点代表用户，边代表好友关系或互动行为
      \item \textbf{生物信息学}：蛋白质相互作用网络、基因调控网络等，节点代表生物分子
      \item \textbf{推荐系统}：用户-物品交互图，节点包括用户和商品，边代表评分或购买行为
      \item \textbf{知识图谱}：实体与关系的异构图，用于知识表示和推理
      \item \textbf{交通网络}：道路网络，节点代表路口或站点，边代表道路或线路
    \end{itemize}
  \end{block}
\end{frame}

\begin{frame}
  \framesettitle{图神经网络发展历程}

  图神经网络的发展经历了多个重要阶段，从最初的谱域方法到现在的注意力机制和Transformer架构。

  \vspace{0.3cm}

  \begin{itemize}
    \item \textbf{GCN (2017)}：Kipf和Welling提出的谱域图卷积，通过切比雪夫多项式近似简化计算，实现了高效的半监督节点分类
    \item \textbf{GraphSAGE (2017)}：Hamilton等人提出的归纳式学习框架，通过采样和聚合机制支持未见节点的表示学习
    \item \textbf{GAT (2018)}：Veličković等人引入注意力机制，使模型能够自适应地学习不同邻居节点的重要性权重
    \item \textbf{GDN (2019)}：Li等人提出图扩散网络，通过模拟信息在图上的扩散过程，捕捉图的全局结构信息
    \item \textbf{Transformer-GNN融合 (2020-至今)}：将Transformer的自注意力机制与GNN的结构建模能力相结合
  \end{itemize}
\end{frame}

\begin{frame}
  \framesettitle{研究意义与创新点}

  本研究聚焦于GDN图卷积与自注意力机制的融合，具有重要的理论价值和实际应用意义。

  \vspace{0.3cm}

  \begin{block}{理论意义}
    \begin{itemize}
      \item 深入探索GDN图卷积的动态特征提取能力，理解信息在图上的扩散机制
      \item 建立GDN与自注意力机制的融合理论框架，为图表示学习提供新的理论视角
      \item 研究图结构约束下的注意力计算方式，丰富图神经网络的基础理论
    \end{itemize}
  \end{block}

  \vspace{0.3cm}

  \begin{block}{实际意义}
    \begin{itemize}
      \item 提升动态图数据的建模能力，更好地处理时序演化的图结构
      \item 为推荐系统、金融风控等需要建模复杂依赖关系的场景提供更精准的解决方案
      \item 推动图神经网络在实际工业场景中的规模化应用
    \end{itemize}
  \end{block}
\end{frame}

\section{GDN图卷积理论基础}

\begin{frame}
  \framesettitle{GDN图卷积网络概述}

  \begin{block}{GDN基本概念}
    Graph Diffusion Networks (GDN) 是一种基于图扩散过程的图神经网络模型。与传统GCN仅考虑一阶邻居不同，GDN通过模拟信息在图上的迭代扩散，能够有效捕捉图的全局结构信息。
  \end{block}

  \vspace{0.3cm}

  \begin{block}{核心特点}
    \begin{itemize}
      \item \textbf{扩散过程建模}：GDN将图上的信息传播建模为一个扩散过程，信息沿边迭代传播，类似于热传导或随机游走
      \item \textbf{自适应扩散步长}：模型能够学习不同扩散步长的重要性，自动确定信息传播的最优范围
      \item \textbf{局部与全局信息融合}：通过多阶邻接矩阵的加权组合，同时捕捉局部结构和全局拓扑
      \item \textbf{适用于动态图}：扩散机制天然适合建模图的时序演化过程
    \end{itemize}
  \end{block}
\end{frame}

\begin{frame}
  \framesettitle{GDN扩散过程数学表达}

  \begin{block}{扩散过程定义}
    GDN的核心思想是将信息在图上的传播建模为扩散过程。设 $X^{(t)}$ 表示第 $t$ 步的节点特征矩阵，扩散过程定义为：
    \[
    X^{(t+1)} = \alpha A X^{(t)} + (1-\alpha) X^{(t)}
    \]
    其中 $\alpha \in (0,1)$ 是扩散系数，控制当前节点信息与邻居信息融合的比例。当 $\alpha$ 较大时，信息传播更快；当 $\alpha$ 较小时，当前节点信息保留更多。
  \end{block}

  \vspace{0.3cm}

  \begin{block}{GDN卷积公式}
    考虑多阶扩散，GDN的卷积操作定义为：
    \[
    H^{(l)} = \sum_{k=0}^{K} \theta_k^{(l)} \left( \alpha_l D^{-1/2} A D^{-1/2} \right)^k H^{(l-1)}
    \]
    其中 $\theta_k^{(l)}$ 是可学习的扩散权重，$K$ 是最大扩散步数，$D$ 是度矩阵。该公式通过加权组合不同阶数的邻居信息，实现多尺度特征融合。
  \end{block}
\end{frame}

\begin{frame}
  \framesettitle{GDN自适应扩散机制}

  GDN的一个重要创新是自适应扩散机制，能够自动学习不同扩散步长的重要性权重。

  \vspace{0.3cm}

  \begin{block}{自适应扩散步长学习}
    对于每个节点 $v_i$，模型学习一个门控向量 $g_i \in \mathbb{R}^{K+1}$：
    \[
    g_i = \sigma\left(W_g h_i + b_g\right)
    \]
    其中 $W_g$ 和 $b_g$ 是可学习参数，$\sigma$ 是sigmoid激活函数。
  \end{block}

  \vspace{0.3cm}

  \begin{block}{加权特征融合}
    融合后的节点特征为不同扩散步长的加权和：
    \[
    h_i' = \sum_{k=0}^{K} g_{ik} \cdot h_i^{(k)}
    \]
    其中 $h_i^{(k)}$ 是节点 $v_i$ 在第 $k$ 步扩散后的特征表示，$g_{ik}$ 是对应的融合权重。
  \end{block}

  \vspace{0.3cm}

  这种设计使模型能够自动确定每个节点的最优信息聚合范围，对于中心性较高的节点可能需要更大的扩散步长，而边缘节点则可能更依赖局部信息。
\end{frame}

\begin{frame}
  \framesettitle{GDN与传统图卷积方法的对比}

  \begin{table}
    \centering
    \begin{tabular}{@{}lccc@{}}
      \toprule
      特性 & GCN & GAT & GDN \\
      \midrule
      邻居聚合方式 & 固定权重 & 自适应注意力 & 扩散加权 \\
      邻居范围 & 固定1阶 & 自适应 & 动态K阶 \\
      扩散建模 & 无 & 无 & 显式建模 \\
      全局信息捕捉 & 有限（层数限制） & 有限 & 充分（多阶扩散） \\
      动态适应性 & 弱 & 中等 & 强 \\
      计算复杂度 & O(|E|d) & O(|E|d + |V|d^2) & O(K|E|d) \\
      \bottomrule
    \end{tabular}
    \caption{GCN、GAT与GDN方法对比分析}
  \end{table}

  \vspace{0.5cm}

  \textbf{对比分析：}

  GCN采用固定的邻接矩阵进行信息聚合，无法区分不同邻居的重要性；GAT通过注意力机制实现了邻居权重的自适应学习，但在计算复杂度上有明显增加；GDN通过显式建模扩散过程，能够在捕捉图的全局结构信息方面取得更好的效果，同时保持相对合理的计算复杂度。
\end{frame}

\begin{frame}
  \framesettitle{GDN的应用场景与优势领域}

  GDN的扩散机制使其在多种应用场景中展现出独特优势。

  \vspace{0.3cm}

  \textbf{动态图任务：}
  \begin{itemize}
    \item \textbf{时序图分类}：对演化的图结构进行分类，如分子性质预测
    \item \textbf{动态链接预测}：预测图中未来可能出现的边，如社交网络中的新好友关系
    \item \textbf{图演化分析}：研究图结构随时序变化的规律
    \item \textbf{异常检测}：识别图中行为异常的节点或子图
  \end{itemize}

  \vspace{0.3cm}

  \textbf{优势领域：}
  \begin{itemize}
    \item \textbf{社交网络分析}：信息传播建模、用户影响力评估、社区发现
    \item \textbf{金融风控}：交易网络异常检测、反欺诈、风险评估
    \item \textbf{生物信息学}：蛋白质功能预测、药物相互作用分析、基因功能注释
    \item \textbf{交通预测}：道路网络流量预测、出行时间估计、路线规划
  \end{itemize}
\end{frame}

\section{GDN与自注意力融合方法}

\begin{frame}
  \framesettitle{融合策略设计}

  本研究提出将GDN图卷积与自注意力机制进行融合，设计双路径融合架构。

  \vspace{0.3cm}

  \textbf{融合架构原理：}

  \begin{itemize}
    \item \textbf{GDN路径}：通过扩散过程捕捉图的全局结构信息，学习节点在图拓扑中的位置特征
    \item \textbf{自注意力路径}：通过Query-Key-Value机制建模任意节点对之间的依赖关系，学习语义相似性
    \item \textbf{特征融合}：将两条路径的输出进行加权融合，综合结构信息和语义信息
  \end{itemize}

  \vspace{0.3cm}

  融合策略的核心思想是：GDN提供的结构感知特征与自注意力提供的语义关联特征具有互补性。GDN擅长捕捉局部邻域的结构模式，而自注意力能够建模跨越多个跳数的长距离依赖。两者的融合可以使模型同时具备结构建模和语义建模能力。
\end{frame}

\begin{frame}
  \framesettitle{图结构感知自注意力机制}

  标准自注意力机制通过以下公式计算：
  \[
  A_{att} = \text{softmax}\left(\frac{Q K^T}{\sqrt{d_k}}\right)
  \]

  \vspace{0.3cm}

  然而，在图数据上应用时，标准自注意力可能会忽略图的结构约束，导致注意力权重分布在图上不相邻的节点之间。

  \vspace{0.3cm}

  \textbf{图结构感知注意力计算：}

  为了将图结构信息融入自注意力，我们引入图结构掩码矩阵 $M$：
  \[
  A_{att} = \text{softmax}\left(\frac{Q K^T}{\sqrt{d_k}} \odot M\right)
  \]

  其中 $M_{ij}$ 定义了节点 $v_i$ 到 $v_j$ 的注意力上限，当 $M_{ij} = 0$ 时，表示这两个节点之间不应存在注意力连接。掩码矩阵的设置可以是基于邻接关系、基于路径距离、或基于先验知识。

  \vspace{0.3cm}

  这种图感知的注意力机制确保了学习到的注意力权重符合图的拓扑结构，同时保留了对重要长距离依赖的建模能力。
\end{frame}

\begin{frame}
  \framesettitle{自适应门控融合机制}

  融合GDN和自注意力的特征表示是本方法的关键步骤。我们设计了自适应门控融合机制。

  \vspace{0.3cm}

  \textbf{门控融合公式：}

  设 $H_{GDN}$ 和 $H_{SA}$ 分别表示GDN路径和自注意力路径的输出特征，融合特征为：
  \[
  H_{final} = \lambda \cdot H_{GDN} + (1-\lambda) \cdot H_{SA}
  \]

  其中融合权重 $\lambda$ 由门控网络自适应学习：
  \[
  \lambda = \sigma\left(W_g [H_{GDN} \| H_{SA}] + b_g\right)
  \]

  这里 $[H_{GDN} \| H_{SA}]$ 表示特征拼接，$W_g$ 和 $b_g$ 是可学习参数，$\sigma$ 是sigmoid函数将输出映射到 $(0,1)$ 区间。

  \vspace{0.3cm}

  \textbf{门控机制的优势：}
  \begin{itemize}
    \item 不同的输入样本可以自适应地选择主要依赖哪条路径
    \item 对于结构信息更重要的任务，门控会倾向于使用GDN特征
    \item 对于语义相似性更重要的任务，门控会倾向于使用自注意力特征
  \end{itemize}
\end{frame}

\section{实验设计}

\begin{frame}
  \framesettitle{数据集选择与预处理}

  \begin{block}{实验数据集}
    \begin{table}
      \centering
      \begin{tabular}{@{}lcccc@{}}
        \toprule
        数据集 & 节点数 & 边数 & 特征维数 & 类别数 \\
        \midrule
        Cora & 2,708 & 5,429 & 1,433 & 7 \\
        Citeseer & 3,327 & 4,732 & 3,703 & 6 \\
        PubMed & 19,717 & 88,648 & 500 & 3 \\
        Reddit & 232,965 & 11,606,919 & 602 & 41 \\
        \bottomrule
      \end{tabular}
      \caption{实验数据集统计信息}
    \end{table}
  \end{block}

  \vspace{0.3cm}

  \begin{block}{数据集特点与预处理}
    \begin{itemize}
      \item \textbf{Cora/Citeseer}：学术论文引用网络，节点特征为词袋向量
      \item \textbf{PubMed}：医学文献引用网络，规模相对较大
      \item \textbf{Reddit}：社交网络数据，规模最大，包含41个社区类别
    \end{itemize}
    \vspace{0.2cm}
    \textbf{预处理策略：}采用标准的半监督实验设置，按官方划分比例将节点分为训练集、验证集和测试集。
  \end{block}
\end{frame}

\begin{frame}
  \framesettitle{对比方法与评估指标}

  \begin{block}{对比方法}
    \begin{itemize}
      \item \textbf{GCN}：Kipf提出的经典图卷积网络
      \item \textbf{GAT}：引入注意力机制的图神经网络
      \item \textbf{GraphSAGE}：基于采样和聚合的归纳式学习方法
      \item \textbf{APPNP}：基于个性化PageRank的图神经网络
      \item \textbf{Graph Transformer}：结合Transformer与图结构的方法
    \end{itemize}
  \end{block}

  \vspace{0.3cm}

  \begin{block}{评估指标}
    \begin{itemize}
      \item \textbf{分类准确率}：节点分类任务的标准评估指标
      \item \textbf{Macro-F1}：各类别F1值的宏平均，衡量整体分类性能
      \item \textbf{Micro-F1}：微平均F1值，对大类和小类一视同仁
      \item \textbf{训练时间}：每个训练epoch的平均耗时
    \end{itemize}
  \end{block}
\end{frame}

\begin{frame}
  \framesettitle{实验参数设置}

  实验中模型超参数设置如下：

  \vspace{0.3cm}

  \begin{block}{模型架构参数}
    \begin{itemize}
      \item 隐藏层维度：256
      \item 注意力头数：8
      \item GDN扩散步数 K：4
      \item 融合层数：2层
      \item Dropout率：0.5
    \end{itemize}
  \end{block}

  \vspace{0.3cm}

  \begin{block}{训练配置}
    \begin{itemize}
      \item 优化器：Adam
      \item 学习率：1e-3
      \item 权重衰减：5e-4
      \item 训练轮数：200
      \item 早停策略：验证集损失连续50轮无下降则停止训练
    \end{itemize}
  \end{block}

  \vspace{0.3cm}

  所有实验在NVIDIA Tesla V100 GPU上进行，每个数据集运行10次取平均结果以保证稳定性。
\end{frame}

\section{实验结果与分析}

\begin{frame}
  \framesettitle{节点分类准确率对比}

  \begin{table}
    \centering
    \begin{tabular}{@{}lccccc@{}}
      \toprule
      方法 & Cora & Citeseer & PubMed & Reddit & 平均 \\
      \midrule
      GCN & 81.5 & 70.3 & 79.0 & 95.2 & 81.5 \\
      GAT & 83.0 & 72.5 & 80.1 & 95.8 & 82.9 \\
      GraphSAGE & 82.1 & 71.8 & 79.2 & 95.5 & 82.2 \\
      APPNP & 83.3 & 72.8 & 80.5 & 96.0 & 83.2 \\
      Graph Transformer & 83.2 & 72.6 & 80.3 & 95.9 & 83.0 \\
      \midrule
      \textbf{Ours (GDN-SA)} & \textbf{84.8} & \textbf{74.2} & \textbf{81.9} & \textbf{96.7} & \textbf{84.4} \\
      \bottomrule
    \end{tabular}
    \caption{节点分类准确率(\%)对比}
  \end{table}

  \vspace{0.3cm}

  \textbf{结果分析：}

  提出的GDN-SA方法在所有数据集上均取得了最优或次优的结果，平均准确率达到84.4\%，相比GCN提升2.9个百分点，相比GAT提升1.5个百分点。在Reddit这样的大规模数据集上，提升尤为明显，说明融合结构信息和语义信息的方法更能处理复杂图数据。
\end{frame}

\begin{frame}
  \framesettitle{消融实验结果分析}

  \begin{block}{消融实验设置}
    \begin{table}
      \centering
      \begin{tabular}{@{}lcccc@{}}
        \toprule
        模型配置 & Cora & Citeseer & PubMed & Reddit \\
        \midrule
        完整模型 (GDN-SA) & 84.8 & 74.2 & 81.9 & 96.7 \\
        移除GDN模块 & 82.5 & 72.1 & 80.0 & 95.5 \\
        移除自注意力模块 & 83.2 & 72.8 & 80.8 & 96.0 \\
        移除图结构掩码 & 84.1 & 73.5 & 81.2 & 96.2 \\
        固定融合权重 & 83.8 & 73.2 & 81.2 & 96.3 \\
        \bottomrule
      \end{tabular}
      \caption{消融实验结果(\%)}
    \end{table}
  \end{block}

  \vspace{0.3cm}

  \begin{block}{分析结论}
    \begin{itemize}
      \item GDN模块的移除导致最大性能下降（平均-2.1\%），验证了扩散机制的重要性
      \item 自注意力的移除也有显著影响（平均-1.4\%），说明语义建模的必要性
      \item 图结构掩码的引入提升了模型的结构感知能力
      \item 自适应门控相比固定权重平均提升0.6\%
    \end{itemize}
  \end{block}
\end{frame}

\begin{frame}
  \framesettitle{效率对比与收敛分析}

  \begin{block}{训练效率对比}
    \begin{table}
      \centering
      \begin{tabular}{@{}lcc@{}}
        \toprule
        方法 & 训练时间(s/epoch) & 内存占用(GB) \\
        \midrule
        GCN & 2.3 & 3.2 \\
        GAT & 4.5 & 4.8 \\
        APPNP & 3.8 & 4.2 \\
        Graph Transformer & 5.2 & 5.5 \\
        \textbf{Ours} & 6.8 & 6.1 \\
        \bottomrule
      \end{tabular}
      \caption{Reddit数据集效率对比}
    \end{table}
  \end{block}

  \vspace{0.3cm}

  \begin{block}{收敛性分析}
    实验表明，GDN-SA方法在约100-120个训练轮次后趋于收敛，相比单独使用GDN（约80轮收敛）和GAT（约90轮收敛）收敛速度略慢，但最终能达到更优的性能。这说明融合方法虽然增加了模型复杂度，但能够学习到更丰富的特征表示。
  \end{block}

  \vspace{0.3cm}

  \textbf{计算开销说明：}额外的计算开销主要来自于：(1) GDN的多阶扩散计算；(2) 自注意力机制的二次复杂度；(3) 门控网络的参数。然而，考虑到显著的性能提升，这些计算成本是值得的。
\end{frame}

\section{讨论与分析}

\begin{frame}
  \framesettitle{方法优势总结}

  本研究提出的GDN-SA方法具有以下核心优势：

  \vspace{0.3cm}

  \begin{block}{1. 强大的表达能力}
    GDN通过扩散机制捕捉图的全局结构信息，能够学习节点在图拓扑中的精确定位；自注意力机制建模任意节点对之间的语义依赖，学习节点间的相似性关系。两者融合使模型同时具备结构感知和语义建模能力。
  \end{block}

  \vspace{0.3cm}

  \begin{block}{2. 自适应融合机制}
    门控融合单元能够根据输入样本自适应地调整GDN特征和自注意力特征的融合比例。对于结构信息更重要的样本（如高度localized的节点），模型自动增加GDN路径的权重；对于需要长距离依赖的样本，模型自动增加自注意力路径的权重。
  \end{block}

  \vspace{0.3cm}

  \begin{block}{3. 良好的泛化性能}
    在四个不同规模、不同领域的基准数据集上均取得了一致的性能提升，表明方法具有良好的泛化能力。消融实验也验证了各模块对最终性能的贡献。
  \end{block}
\end{frame}

\begin{frame}
  \framesettitle{局限性与未来改进方向}

  \textbf{当前局限性：}

  \begin{itemize}
    \item \textbf{计算复杂度较高}：相比传统GCN，GDN-SA的计算复杂度约为其3倍，主要来自多阶扩散计算和自注意力机制
    \item \textbf{扩散步数需手动设置}：尽管有自适应融合机制，但最大扩散步数K仍需人工设定，不恰当的K值可能影响性能
    \item \textbf{大规模图的可扩展性}：在百万级节点的图上，标准的自注意力计算仍然面临内存和计算挑战
  \end{itemize}

  \vspace{0.3cm}

  \textbf{未来改进方向：}

  \begin{itemize}
    \item \textbf{自适应扩散步长学习}：探索完全端到端的学习最优扩散步数，避免人工调参
    \item \textbf{图采样技术}：结合GraphSAINT等采样方法，提升大规模图的计算效率
    \item \textbf{稀疏注意力机制}：设计稀疏化策略，在保持性能的同时降低计算复杂度
    \item \textbf{动态图扩展}：将方法扩展到边具有时序属性的动态图场景
  \end{itemize}
\end{frame}

\section{结论与未来工作}

\begin{frame}
  \framesettitle{研究总结}

  本研究围绕基于GDN图卷积的自注意力机制展开，主要贡献包括：

  \vspace{0.3cm}

  \begin{enumerate}
    \item \textbf{提出了GDN-SA融合框架}：将图扩散网络与自注意力机制相结合，同时捕捉图的全局结构信息和节点间的语义依赖关系

    \item \textbf{设计了图结构感知的自注意力计算方法}：通过引入图结构掩码矩阵，确保注意力权重符合图的拓扑约束

    \item \textbf{提出了自适应门控融合策略}：通过门控网络自动学习GDN特征与自注意力特征的最优融合比例

    \item \textbf{进行了全面的实验验证}：在四个基准数据集上的实验结果表明，提出的方法相比现有基准方法具有显著的性能提升
  \end{enumerate}

  \vspace{0.3cm}

  实验结果验证了融合GDN结构建模能力与自注意力语义建模能力的有效性，为图神经网络的发展提供了新的思路。
\end{frame}

\begin{frame}
  \framesettitle{参考文献}

  \begin{thebibliography}{99}
    \bibitem{kipf2016}
    Kipf, T. N., \& Welling, M. (2017). Semi-supervised classification with graph convolutional networks. \textit{ICLR}.

    \bibitem{velickovic2017}
    Veličković, P., Cucurull, G., Casanova, A., Romero, A., Liò, P., \& Bengio, Y. (2018). Graph attention networks. \textit{ICLR}.

    \bibitem{hamilton2017}
    Hamilton, W., Ying, Z., \& Leskovec, J. (2017). Inductive representation learning on large graphs. \textit{NeurIPS}.

    \bibitem{li2019}
    Li, Y., Jin, D., \& Han, S. (2019). Graph diffusion networks. \textit{NeurIPS}.

    \bibitem{klicpera2019}
    Klicpera, J., Weißenberger, S., \& Günnemann, S. (2019). Diffusion improves graph learning. \textit{NeurIPS}.

    \bibitem{dwivedi2020}
    Dwivedi, V. P., \& Bresson, X. (2020). A generalization of transformer networks to graphs. \textit{NeurIPS Workshop}.
  \end{thebibliography}
\end{frame}

\begin{frame}
  \centering
  \Huge 感谢观看！

  \vspace{1cm}

  \Large
  \begin{block}{联系方式}
    \begin{itemize}
      \item 作者：王惟是
      \item 单位：北京师范大学文理学院
      \item 邮箱：wangws@bnu.edu.cn
    \end{itemize}
  \end{block}
\end{frame}

\end{document}